\documentclass[11pt]{article}

\usepackage[utf8]{inputenc}
\usepackage[T1]{fontenc}

\usepackage[english]{babel}
\usepackage{indentfirst}

\usepackage[
	a4paper,
	left = 2cm,
	right = 2cm,
	bottom = 1cm, includefoot,
	top = 1.2cm, includehead,
]{geometry}

\usepackage{amsmath,amssymb,amsfonts}
\RequirePackage{physics}
\usepackage{siunitx}
\usepackage[version=4]{mhchem}
\usepackage{chemfig}
\AtBeginDocument{\RenewCommandCopy\qty\SI}
\AtBeginDocument{\RenewCommandCopy\ce\chem}

\usepackage{enumitem}
\setlist{
	leftmargin = *,
}
\numberwithin{enumi}{section}

\usepackage{booktabs}

\usepackage{csquotes}
\usepackage{hyperref}

\title{\textbf{Optical Elements - Prisms}}
\author{Rodrigo Sánchez-Martínez}
\date{\today}

% --
\begin{document}
	
\maketitle	
	
%.
\tableofcontents

\section*{Introduction}

	\begin{itemize}
	
	\item 
	Prisms are solid glass optics that are ground and polished into geometrical and optically significant shapes.
	The angle, position, and number of surfaces help define the type and function.
	They can be found in various shapes and configurations, such as right-angle, dispersing, roof, dove, retro-reflecting, penta, and wedge prisms as well as hexagonal light mixing rods.
	
	\item
	One of the most recognizable uses of prisms, as demonstrated by Sir Isaac Newton, consists of dispersing a beam of white light into its component colors.
	This application is utilized by refractometer and spectrographic components.
	
	\item
	Since this initial discovery, prisms have been used in \enquote{bending} light within a system, \enquote{folding} the system into a smaller space, changing the orientation (also known as handedness or parity) of an image, as well as combining or splitting optical beams with partial reflecting surfaces.
	These uses are common in applications with telescopes, binoculars, surveying equipment, and a host of others.

	\end{itemize}
	
\section{}

\end{document}